%-----------------------------------------------------------------------------------------------------
\chapter{Conclusion and Future Work}
\label{cap:conclusao}
%-----------------------------------------------------------------------------------------------------

%-----------------------------------------------------------------------------------------------------
\section{Restatement of Objectives}
%-----------------------------------------------------------------------------------------------------

The primary objective of this work was to develop a mechatronic and computerised system for the measurement of antenna radiation patterns. The system was conceived with a focus on low cost, aiming to provide an accessible and practical tool for academic environments, particularly in courses related to telecommunications, radio frequency, and antenna theory. By doing so, the work saw the opportunity to bridge the gap between the high cost of commercial equipment and the educational need for hands-on demonstrations of radiation patterns, which are often difficult to understand only in theoretical terms.

%-----------------------------------------------------------------------------------------------------
\section{Summary of Achievements}
%-----------------------------------------------------------------------------------------------------

The proposed system was thought, designed, implemented, and tested, providing experimental radiation pattern measurements of an antenna operating at the ISM frequency of 868~MHz. The results obtained showed good agreement with simulated data, acknowledging the limitations of measurements carried out in an outdoor, non-anechoic environment.

The incremental characterisation of the system components provided a clear understanding of their roles:
\begin{itemize}
    \item The antenna alone demonstrated good matching and radiation characteristics close to the intended frequency.
    \item The band-pass filter introduced the expected insertion loss while improving frequency selectivity.
    \item The balun exhibited low differential-mode loss and strong common-mode suppression, enhancing the current distribution.
    \item The impedance matching network improved the alignment with the 50~$\Omega$ reference, optimising power transfer in the measurement setup.
\end{itemize}

Altogether, the system proved to be effective in reproducing the main radiation features while remaining affordable and easy to deploy, thus offering significant pedagogical potential.

%-----------------------------------------------------------------------------------------------------
\section{Limitations}
%-----------------------------------------------------------------------------------------------------

Although the objectives were met, certain limitations must be acknowledged:
\begin{itemize}
    \item The mechanical structure does not easily accommodate larger antennas, which restricts its versatility.
    \item The measurement process, while sufficiently accurate, requires a relatively long acquisition time due to the stepwise mechanical actuation.
    \item The usable frequency range is constrained by the selected components and available equipment.
    \item The absence of an anechoic chamber meant that measurements were subject to multipath and environmental reflections.
\end{itemize}

These factors do not diminish the value of the system for its intended purpose but provide guidance for future improvements.

%-----------------------------------------------------------------------------------------------------
\section{Future Work}
%-----------------------------------------------------------------------------------------------------
Several avenues for future development and research are suggested:
\begin{itemize}
    \item Performing measurements in a controlled environment, such as an anechoic chamber, in order to obtain more accurate radiation patterns.
    \item Extending the system to support a broader range of frequencies and different antenna types.
    \item Reducing the acquisition time by enhancing the speed and automation of the mechanical rotation system, while maintaining measurement accuracy.
    \item Integration of a low-cost RF signal generator operating at 868~MHz, with the aim of further reducing the dependence on expensive laboratory instruments.
    \item Developing a user-friendly software interface tailored for educational use, enabling integration into laboratory classes.
\end{itemize}

%-----------------------------------------------------------------------------------------------------
\section{Closing Remarks}
%-----------------------------------------------------------------------------------------------------

In summary, this work demonstrated the viability of a low-cost, mechatronic antenna measurement system that can support educational activities in RF and telecommunications. The developed system fills an important gap by providing an accessible and practical tool to visualise antenna radiation patterns, fostering greater understanding of a concept that is otherwise abstract and difficult to convey.

Beyond the technical results, the project also contributed to the professional development of engineering competencies. The process involved the use of CAD tools such as SolidWorks, electromagnetic simulation software such as CST Studio Suite, and practical experience with measurement instruments including VNAs, SAs, and many others. These experiences are of significant value to the formation of an engineer and further strengthen the educational impact of the work.

In addition, the project provided valuable learning experiences regarding the practical behaviour of RF components. Unlike in DC or low-frequency circuits, where idealised behaviour is often assumed, at higher frequencies parasitic effects and impedance mismatches play a crucial role. Gaining hands-on familiarity with these aspects reinforced the theoretical concepts learned in the classroom and highlighted the importance of careful design and characterisation in RF systems.

The outcomes presented here build a solid foundation for further improvement of the system and open opportunities for different applications in teaching and research, reinforcing the value of accessible tools for engineering education.

